% !TEX program = xelatex
% Алгоритм оптимизации трёхмерного расположения зеркальной системы LiDAR.
% Сборка: xelatex -output-directory=build article.tex
%         (два прохода)

\documentclass{misisarticle}

\grnti{28.23.27, 28.23.15}
\asjc{1707, 1711}
\oecd{1.2}

\title{Алгоритм оптимизации трехмерного расположения зеркальной системы LiDAR для детализации элементов сцены}

\misisauthor{Д.~В.~Плишко}{кафедра инженерной кибернетики НИТУ «МИСиС»}{Москва, Россия}{m2304259@edu.misis.ru}

\abstract{%
  Рассматривается автоматический подбор расположения и ориентации нескольких LiDAR и плоских зеркал, чтобы полнее детализировать элементы трёхмерной сцены. Часть приборов неподвижна, часть движется вместе с носителем, как лидары и боковые зеркала автомобиля. Свободные позы выбираются под конкретную геометрию, в одном положении или вдоль траектории. Покрытие тем лучше, чем выше доля наблюдаемых поверхностных вокселей и равномерность лучей, чем меньше площадь слепых зон и чем больше зеркала добавляют к прямому сканированию. Сцена задаётся треугольным мешем. Видимость луча определяется ближайшим пересечением с мешем и законом отражения на зеркалах. Грубое размещение строится жадным покрытием, непрерывное уточнение выполняется имитацией отжига или CMA-ES. Пересечение с мешем на центральном процессоре считает Intel Embree. Тот же запрос предусматривается на графическом процессоре: Vulkan Ray Tracing для трассировки и Vulkan Compute для свёртки покрытия.
}

\keywords{%
  LiDAR, Intel Embree, BVH, Vulkan, воксель, имитация отжига
}

\usepackage{tikz}
\usetikzlibrary{arrows.meta,positioning}

\begin{document}

\maketitle

% ============================================================
\section{Введение}

Лидар детализирует только ту часть сцены, куда доходит прямой луч. Поле зрения ограничено, а крупные элементы закрывают пространство за собой. Ещё один сенсор закрывает часть таких зон, но занимает место, питание и отдельный канал. Плоское зеркало оставляет уже имеющиеся лучи и уводит их туда, куда прямое сканирование не смотрит.

На носителе таких приборов обычно несколько. У автомобиля это лидары на крыше и по бортам и штатные боковые зеркала: часть из них жёстко связана с кузовом и меняет положение вместе с ним. Поза, удачная для одной геометрии и одного кадра, на другой сцене или дальше по траектории смотрит в стену либо в уже видимый фасад. Расставлять приборы вручную для каждого случая нельзя. В работе формулируется подбор свободных поз набора лидаров и зеркал по покрытию трёхмерной сцены и описывается метод: поверхностные воксели, трассировка с отражением и двухфазный поиск.

% ============================================================
\section{Анализ работ}

Размещение лидара и применение зеркала пока решаются как разные задачи. Ближние работы сведены в таблице~\ref{tab:related}. Ни одна из них не ищет позу физического зеркала по покрытию элементов произвольной трёхмерной сцены.

\begin{table}[htbp]
  \caption{Ближние постановки}
  \label{tab:related}
  \centering
  \footnotesize
  \setlength{\tabcolsep}{3pt}
  \begin{tabular}{@{}lll@{}}
    \toprule
    \textbf{Работа} & \textbf{Что размещают} & \textbf{Оценка} \\
    \toprule
    \cite{vishnevsky2026} & виртуальный канал & разрешение \\
    \cite{li2021} & зеркала сканера & охват плоскости \\
    \cite{cai2023} & LiDAR инфраструктуры & точность детектора \\
    \cite{kini2020} & несколько LiDAR & занятость облака \\
    эта работа & лидары и зеркала & воксели сцены \\
    \bottomrule
  \end{tabular}
\end{table}

Вишневский и Фёдоров повышают пространственное разрешение механического LiDAR виртуальным зеркальным каналом и автоматической калибровкой~\cite{vishnevsky2026}. Канал виртуальный: он уплотняет отсчёты уже существующей развёртки, а не выводит луч в закрытую область. Калибровка восстанавливает параметры канала. Положение зеркала по метрике слепых зон сцены не ищется.

Ли, Ли, Ким и Ло измеряют плоскостность лазерным сканером и ставят плоские зеркала, чтобы увеличить дальность и охват~\cite{li2021}. Зеркало здесь физическое и действительно добавляет лучи туда, куда прямая видимость не достаёт. Раскладка подчинена контролю одной поверхности. Оптимизатора позы под произвольную геометрию и отдельного выигрыша зеркала в этой работе нет.

Цай и соавторы ищут места для инфраструктурных LiDAR~\cite{cai2023}. Для оценки построена библиотека имитации разных лидаров в симуляторе CARLA. Качество схемы измеряется точностью детекторов, а связь с облаком точек~--- плотностью и равномерностью области интереса. На типовой полосе, при том же числе и типе сенсоров, подобранная схема подняла среднюю точность примерно на 15\% относительно обычной расстановки. Объект другой: неподвижные сенсоры инфраструктуры и точность восприятия. Отражения и зеркала в постановке нет. Плотность и равномерность как признаки качества облака к нашей оценке применимы.

Кини решает, куда ставить один и несколько LiDAR на транспортном средстве, тоже в CARLA~\cite{kini2020}. Целевая функция считается по занятости облака вокруг машины, поиск ведёт генетический алгоритм. Проверялись грузовик и легковой автомобиль, сенсоры Velodyne с 16 и 32 каналами и разные области интереса. Это ближайшая постановка по вопросу «где разместить сенсор». Сенсоров несколько, зеркала нет, а прямой и отражённый вклад не разделяются.

Отсюда рамка данной работы. В системе несколько лидаров и несколько плоских зеркал, и часть из них едет вместе с носителем, а не стоит в одном кадре. Зеркало задаётся конечным прямоугольником и законом отражения, а не виртуальным каналом. Оценка геометрическая и относится к элементам сцены: доля наблюдаемых вокселей, площадь слепых зон, равномерность лучей и отдельный выигрыш зеркал над прямым сканом. Свободных поз уже слишком много для плотной сетки, поэтому размещение ищется в два шага.

% ============================================================
\section{Постановка задачи}

\subsection{Что требуется получить}

Нужно автоматически подобрать расположение и ориентацию лидаров и плоских зеркал так, чтобы прямые и отражённые лучи детализировали как можно большую часть поверхности сцены. Зеркало здесь не замена сенсору. Оно направляет уже имеющиеся лучи туда, куда прямой обзор не достаёт: за стену, в проём, вдоль борта, вне поля зрения.

Задача не сводится к одной паре «сенсор и зеркало». Лидаров может быть несколько, зеркал тоже. Часть приборов стоит в сцене неподвижно. Часть закреплена на носителе и едет вместе с ним. Лидар на крыше или бампере и боковое зеркало автомобиля перемещаются в мире, хотя относительно кузова их поза может не меняться. Поэтому искомое~--- не одна точка, а набор поз. Для неподвижного прибора поза задаётся в сцене. Для прибора на носителе поза задаётся в системе носителя, а мировая поза в каждый момент получается из движения носителя.

Свободны только те позы, которые разрешено менять. Размеры зеркал заданы заранее. Если сенсоры уже закреплены, подбираются зеркала. Если зеркала заданы кузовом, подбираются сенсоры. Если свободно и то и другое, они ищутся вместе: поле зрения и направление отражения зависят друг от друга.

На неподвижной сцене достаточно одного положения носителя. На траектории, например при проезде автомобиля, одного кадра мало. Элемент поверхности считается наблюдаемым вдоль пути, если его достиг хотя бы один засчитанный луч хотя бы в одном положении носителя. Рядом сохраняется мгновенное покрытие в отдельных точках пути, чтобы схема, хорошая в среднем, не оставляла провал на конкретном участке.

\subsection{С чем ведётся расчёт}

Сцена задана треугольным мешем. Если носитель сам закрывает лучи, его геометрия входит в тот же меш и движется вместе с закреплёнными приборами.

Каждый лидар задаётся позицией, ориентацией, диаграммой направлений в системе сенсора, масштабом углов этой диаграммы и максимальной дальностью. По диаграмме из лидара выпускается набор лучей.

Каждое зеркало~--- прямоугольник конечного размера. Его задают центр, ориентация, ширина и высота. Ориентация строит репер: две оси лежат в плоскости, третья есть нормаль. Луч отражается только от лицевой стороны и только внутри прямоугольника.

Луч либо пересекает меш напрямую, либо сначала встречает зеркало и продолжается по закону отражения. Цепочка из нескольких зеркал допускается, пока не исчерпан предел отражений и не израсходована дальность. Попадание в поверхность засчитывается, если путь не длиннее дальности и косинус угла падения не ниже порога. Порог отсекает скользящие лучи, которые плохо детализируют элемент.

Наблюдаемость относится к поверхности, а не к пустому объёму. Область интереса разбита на воксели, и в расчёт входят те, через которые проходит меш. Воксель наблюдаем, если в него попал хотя бы один засчитанный луч любого лидара, прямой или отражённый. Слепая зона есть суммарная площадь поверхности в остальных вокселях. Зеркало недопустимо, если пересекает меш, выходит за пределы сцены или накрывает сенсор. Элементарный случай той же схемы показан на рисунке~\ref{fig:scheme}.

\begin{figure}[htbp]
  \centering
  \begin{tikzpicture}[x=0.155cm, y=0.155cm, line join=round, line cap=round]
    \fill[black!6] (1,21) rectangle (8,29);
    \draw[thick] (0,0) rectangle (40,30);
    \draw[thick] (18,0) -- (18,10);
    \draw[thick] (18,24) -- (18,30);
    \fill[red!75!black] (10,15) circle (0.85);
    \node[font=\scriptsize, anchor=north] at (10,13.5) {LiDAR};
    \draw[blue!70!black] (10,15) -- (38,29);
    \draw[blue!70!black] (10,15) -- (38,4);
    \draw[green!55!black] (10,15) -- (30,16);
    \draw[green!55!black] (30,16) -- (4,26);
    \draw[orange, very thick] (28.9,9.1) -- (31.1,22.9);
    \draw[orange, ->] (30,16) -- (27,16.5);
    \node[font=\scriptsize, align=center] at (4.3,28.1) {слепая\\зона};
    \node[font=\scriptsize, anchor=west] at (26,27.2) {прямые лучи};
    \node[font=\scriptsize, anchor=west] at (32.2,11) {зеркало};
    \node[font=\scriptsize, fill=white, inner sep=1pt] at (14.2,23.5) {проём};
  \end{tikzpicture}
  \caption{Элементарный случай модели: один лидар и одно зеркало. Несколько приборов и движение носителя используют те же лучи и то же отражение. Синие отрезки~--- прямые попадания, зелёные~--- путь через зеркало и проём, стрелка~--- нормаль.}
  \label{fig:scheme}
\end{figure}

% ============================================================
\section{Описание методов}

\subsection{Поверхностные воксели}

Равномерная сетка покрывает интересующую область. Воксель с индексом \((i,j,k)\) занимает ячейку со стороной \(h\), отложенную от минимума области. В покрытие входят не все ячейки объёма, а те, через которые проходит поверхность меша. Для каждого такого вокселя хранится площадь поверхности внутри него. Слепая зона~--- это непокрытая площадь поверхности, а не пустой объём: точки попадания и так лежат на элементах сцены.

Доля наблюдаемых вокселей
\begin{equation}
  C = \frac{|V_{\mathrm{obs}}|}{|V|},
  \label{eq:coverage}
\end{equation}
где \(V_{\mathrm{obs}} \subseteq V\)~--- воксели, в которые пришёл хотя бы один засчитанный луч. Площадь слепых зон \(A_{\mathrm{blind}}\)~--- сумма площадей вокселей из \(V \setminus V_{\mathrm{obs}}\).

\subsection{Отражение и пересечение}

Ориентация зеркала задаёт ортонормированный репер: две оси лежат в плоскости прямоугольника, третья~--- нормаль \(\mathbf{n}\). Луч \(\mathbf{r}(t) = \mathbf{o} + t\mathbf{d}\) пересекает плоскость при
\begin{equation}
  t = \frac{(\mathbf{c} - \mathbf{o}) \cdot \mathbf{n}}{\mathbf{d} \cdot \mathbf{n}},
  \label{eq:plane}
\end{equation}
если знаменатель не близок к нулю и точка попадания лежит внутри прямоугольника. Удар с тыльной стороны не отражает: зеркало одностороннее.

Отражённое направление
\begin{equation}
  \mathbf{d}' = \mathbf{d} - 2(\mathbf{d} \cdot \mathbf{n})\,\mathbf{n}
  \label{eq:reflect}
\end{equation}
нормируется. Новый отрезок начинается чуть впереди точки попадания, чтобы не засчитать то же зеркало повторно, а остаток дальности уменьшается на пройденный путь. Число отражений ограничено. Одного отражения хватает, когда на пути луча одно зеркало. Предел больше единицы нужен для цепочки, например от бокового зеркала к следующему.

Ближайшее пересечение луча с мешем ищет иерархия ограничивающих объёмов. На центральном процессоре её считает Intel Embree~\cite{wald2014}. Зеркала проверяются аналитически и сравниваются с дистанцией до меша: их меньше, чем треугольников. Если какое-то зеркало ближе и удар лицевой, луч отражается по~\eqref{eq:reflect}, и запрос повторяется. Иначе записывается попадание в сцену: точка, нормаль, полная длина пути, косинус падения и индекс зеркала, если отражение было. Лучи всех лидаров собираются в одно покрытие.

Модель луча не привязана к Embree. Нужны запрос ближайшего пересечения с мешем и проверка прямоугольников. Далее тот же запрос предполагается исполнять на графическом процессоре: иерархию меша и ближайшее попадание даёт Vulkan Ray Tracing, а разнесение попаданий по вокселям и подсчёт метрик кандидата~--- Vulkan Compute. Оба пути должны давать одну оценку позы, чтобы поиск можно было сверять на центральном процессоре и ускорять на графическом.

Покрытие копится после попадания, по идентификатору поверхностного вокселя. Пошаговый обход пустых ячеек (ray marching) здесь не используется: детализировать в пустоте нечего, а пересечение с треугольниками уже даёт BVH. Воксельная сетка дискретизирует оценку, а не шаг луча.

\subsection{Метрики покрытия}

Подбор видит одну скалярную оценку. Её слагаемые отвечают разным дефектам скана и сведены в таблице~\ref{tab:metrics}.

\begin{table}[htbp]
  \caption{Метрики покрытия и направление улучшения}
  \label{tab:metrics}
  \centering
  \begin{tabular}{llc}
    \toprule
    \textbf{Метрика} & \textbf{Обозначение} & \textbf{Лучше} \\
    \toprule
    Доля вокселей & \(C\) & больше \\
    Площадь слепых зон & \(A_{\mathrm{blind}}\) & меньше \\
    Угловая равномерность & \(U\) & больше \\
    Выигрыш зеркала & \(G\) & больше \\
    Составная оценка & \(\Phi\) & больше \\
    \bottomrule
  \end{tabular}
\end{table}

Выигрыш зеркала сравнивает покрытие с зеркалом и без него:
\begin{equation}
  G = \frac{A_{\mathrm{obs}}}{A_{\mathrm{direct}}}, \qquad
  G_{\mathrm{norm}} = \min\left(1,\frac{G - 1}{0{,}5}\right).
  \label{eq:gain}
\end{equation}
Нормировка насыщается, когда зеркало добавляет половину новой площади. Дальнейший рост \(G\) не перетягивает оценку на себя.

Равномерность \(U\) считается по угловым бинам диаграммы сенсора. Бины вне заявленного поля зрения в неё не входят. Дополнительно по попаданиям считается коэффициент вариации числа лучей в наблюдаемых вокселях: он показывает, не собралась ли детализация в одном пятне. Плотность и равномерность облака уже использовались как признаки качества размещения~\cite{cai2023}; здесь они входят в \(\Phi\) рядом с выигрышем зеркала, а не вместо геометрического покрытия.

Составная оценка
\begin{equation}
  \Phi = w_{C} C + w_{U} U + w_{G} G_{\mathrm{norm}}
        + w_{A}\left(1 - \frac{A_{\mathrm{blind}}}{A_{\Sigma}}\right),
  \label{eq:phi}
\end{equation}
где \(A_{\Sigma}\)~--- площадь всех целевых вокселей, а веса неотрицательны и в сумме дают единицу. В постановке принимаются \(w_{C} = 0{,}25\), \(w_{U} = 0{,}15\), \(w_{G} = 0{,}35\), \(w_{A} = 0{,}25\): больший вес стоит на новом обзоре от зеркала.

\subsection{Подбор позы}

Каждый свободный прибор добавляет в поиск свою позу: три координаты и два или три угла. Поза прибора на носителе одна на всю траекторию, а не новая на каждом кадре. Уже одно зеркало при восьми узлах на ось даёт \(8^{5} = 32\,768\) трассировок, а несколько лидаров и зеркал делают полную сетку нереализуемой. Грубая фаза поэтому строит каталог, а не сетку.

Зеркало попадает в каталог, если прямоугольник лежит в свободном пространстве, не пересекает меш и не накрывает сенсор. Поза лидара попадает в каталог, если сенсор лежит в отведённой ему области и не оказывается внутри меша. Центры зеркал берутся среди точек, которые хотя бы один лидар видит напрямую: иначе зеркало само стоит в слепой зоне и не получает лучей. Нормали выбираются так, чтобы отражённый пучок смотрел в воксели с нулевым прямым покрытием.

Каждый кандидат после трассировки даёт множество покрытых целевых вокселей. Жадный шаг берёт позу прибора с наибольшим числом ещё не покрытых вокселей и фиксирует её. Шаг повторяется по свободным лидарам и зеркалам, пока прирост не станет ниже порога. Это жадная схема задачи о покрытии множеств~\cite{chvatal1979}. Если свободен один прибор, она совпадает с выбором лучшего элемента каталога. Набор выбранных поз есть старт непрерывной фазы. На траектории одна такая поза в системе носителя оценивается во всех выбранных положениях движения.

Оценка \(\Phi\) кусочно-постоянна: луч либо попадает в воксель, либо нет. Градиент по позе почти всюду нулевой, поэтому уточнение безынградное.

Имитация отжига хранит текущую позу и температуру \(T\)~\cite{kirkpatrick1983}. Соседняя поза получается малым сдвигом координат и углов. Если \(\Delta\Phi = \Phi_{\mathrm{new}} - \Phi \ge 0\), шаг принимается. Ухудшение принимается с вероятностью
\begin{equation}
  P = \exp(\Delta\Phi / T).
  \label{eq:sa}
\end{equation}
Температура снижается геометрически, \(T \leftarrow \alpha T\), при \(\alpha \in (0,1)\). Высокая начальная температура позволяет уйти от позы, которая смотрит в уже покрытый фасад. К концу поиска остаются только улучшения.

CMA-ES решает ту же непрерывную задачу иначе: выборка из нормального распределения, обновление среднего и ковариации по лучшим точкам~\cite{hansen2001}. Ковариация вытягивается вдоль направлений, где \(\Phi\) растёт. Это полезно, когда сдвиг зеркала вдоль стены почти ничего не меняет, а поворот нормали меняет покрытие резко. Генетический алгоритм с вещественным кодированием позы, уже применённый к расстановке нескольких LiDAR~\cite{kini2020}, допустим как третий безынградный вариант. Порядок фаз показан на рисунке~\ref{fig:pipeline}.

\begin{figure}[htbp]
  \centering
  \begin{tikzpicture}[
      node distance=0.28cm,
      box/.style={
        draw, rounded corners=1.2pt, align=center,
        font=\scriptsize, inner sep=3pt, text width=0.86\columnwidth
      },
      arr/.style={-{Stealth[length=1.5mm]}, thin}
    ]
    \node[box] (mesh) {Меш, набор лидаров и зеркал};
    \node[box, below=of mesh] (cat) {Каталог поз в сцене или на носителе};
    \node[box, below=of cat] (trace) {Пересечение: Embree, далее Vulkan};
    \node[box, below=of trace] (score) {Метрики \(C\), \(A_{\mathrm{blind}}\), \(U\), \(G\)};
    \node[box, below=of score] (greedy) {Жадный выбор позы (set cover)};
    \node[box, below=of greedy] (local) {Уточнение: имитация отжига или CMA-ES};
    \draw[arr] (mesh.south) -- (cat.north);
    \draw[arr] (cat.south) -- (trace.north);
    \draw[arr] (trace.south) -- (score.north);
    \draw[arr] (score.south) -- (greedy.north);
    \draw[arr] (greedy.south) -- (local.north);
  \end{tikzpicture}
  \caption{Порядок расчёта: прямая модель покрытия и двухфазный подбор поз.}
  \label{fig:pipeline}
\end{figure}

\begin{thebibliography}{99}

\bibitem{vishnevsky2026} Вишневский М.Р., Фёдоров И.Ю. Повышение пространственного разрешения механических LiDAR-систем с использованием виртуального зеркального канала и автоматической калибровки. \textit{Труды Института системного программирования РАН}, 2026, т.~38, №~3, с.~101--114. \url{https://doi.org/10.15514/ISPRAS-2026-38(3)-38}

\bibitem{li2021} Li F., Li H., Kim M.-K., Lo K.-C. Laser Scanning Based Surface Flatness Measurement Using Flat Mirrors for Enhancing Scan Coverage Range. \textit{Remote Sensing}, 2021, vol.~13, no.~4, article~714. \url{https://doi.org/10.3390/rs13040714}

\bibitem{cai2023} Cai X., Jiang W., Xu R., Zhao W., Ma J., Liu S., Li Y. Analyzing Infrastructure LiDAR Placement with Realistic LiDAR Simulation Library. \textit{Proceedings of the IEEE International Conference on Robotics and Automation (ICRA)}, 2023. \url{https://doi.org/10.1109/ICRA48891.2023.10161027}

\bibitem{kini2020} Kini R.R. Sensor Position Optimization for Multiple LiDARs in Autonomous Vehicles. Master thesis, KTH Royal Institute of Technology, Stockholm, TRITA-EECS-EX 2020:896, 2020.

\bibitem{wald2014} Wald I., Woop S., Benthin C., Johnson G.S., Ernst M. Embree: A Kernel Framework for Efficient CPU Ray Tracing. \textit{ACM Transactions on Graphics}, 2014, vol.~33, no.~4, article~143.

\bibitem{chvatal1979} Chv\'{a}tal V. A Greedy Heuristic for the Set-Covering Problem. \textit{Mathematics of Operations Research}, 1979, vol.~4, no.~3, pp.~233--235.

\bibitem{kirkpatrick1983} Kirkpatrick S., Gelatt C.D., Vecchi M.P. Optimization by Simulated Annealing. \textit{Science}, 1983, vol.~220, no.~4598, pp.~671--680.

\bibitem{hansen2001} Hansen N., Ostermeier A. Completely Derandomized Self-Adaptation in Evolution Strategies. \textit{Evolutionary Computation}, 2001, vol.~9, no.~2, pp.~159--195.

\end{thebibliography}

\end{document}
